Given $\boldsymbol{d}\in\Pi(3g-3+n,n)$ let
\begin{gather}\label{eq:floor:brackets}
\lfloor \tau_{d_1} \cdots \tau_{d_n}\rfloor_{g,n}=
\frac{(6g-5+2n)!!}{(2d_1+1)!!\cdots(2d_n+1)!!}
\cdot\frac{1}{g!\cdot 24^g} .
\end{gather}
By definition of $\epsilon(\boldsymbol{d})$ we have
\begin{gather}\label{eq:def:of:epsilon}
\langle \tau_{d_1} \cdots \tau_{d_n}\rangle_{g,n}
=
\lfloor \tau_{d_1} \cdots \tau_{d_n}\rfloor_{g,n}
\cdot\big(1+\epsilon(\boldsymbol{d})\big) .
\end{gather}

From now on we suppose that $g\ge 1$.

\begin{Lemma}\label{lm:string}
Let $\boldsymbol{d}\in\Pi(3g-2+n,n-k)$ such that
$d_j>0$ for $j=1,\dots,n-k$. We assume that $k\ge 0$ and $n-k>0$.
Define $\delta_{\rm string}(0^{k+1},\boldsymbol{d})$
by equation
\begin{gather}
\big\lfloor \tau_0^{k+1}\tau_{d_1} \cdots \tau_{d_{n-k}}\big\rfloor_{g,n+1}
\cdot\big(1+\delta_{\rm string}\big(0^{k+1},\boldsymbol{d}\big)\big)
\nonumber\\
\qquad{} =
\big\lfloor \tau_0^k\tau_{d_1-1} \cdots \tau_{d_{n-k}}\big\rfloor_{g,n}
+\dots +\big\lfloor \tau_0^k\tau_{d_1} \cdots \tau_{d_{n-k}-1}\big\rfloor_{g,n} .\label{eq:string:floor}
\end{gather}
Then
\[
\delta_{\rm string}\big(0^{k+1},\boldsymbol{d}\big)=\frac{n-k-1}{6g-3+2n} .
\]
In particular, for any $\boldsymbol{d}$ as above and for any
$k\ge 0$ we have
\begin{gather}\label{eq:delta:string:is:positive}
\delta_{\rm string}\big(0^{k+1},\boldsymbol{d}\big)\ge 0 .
\end{gather}
\end{Lemma}

\begin{proof}Dividing both sides of equation~\eqref{eq:string:floor} by
$\big\lfloor \tau_0^{k+1}\tau_{d_1} \dots \tau_{d_{n-k}}\big\rfloor_{g,n+1}$
and applying definition~\eqref{eq:floor:brackets}
to all terms involved in the right-hand side of the resulting
equation we get
\begin{gather*}
1+\delta_{\rm string}\big(0^{k+1},\boldsymbol{d}\big)
=\frac{(2d_1+1)+\dots+(2d_{n-k}+1)}{6g-3+2n}
 \\
 \hphantom{1+\delta_{\rm string}\big(0^{k+1},\boldsymbol{d}\big)}{} =
\frac{2(3g-2+n)+(n-k)}{6g-3+2n}
=\frac{6g-4+3n-k}{6g-3+2n}
=1+\frac{n-k-1}{6g-3+2n} .\!\!\!\tag*{\qed}
\end{gather*}\renewcommand{\qed}{}
\end{proof}

\begin{Corollary}\label{cor:epsilon:0:2corr}
For any $(d_1,d_2)\in\Pi(3g-2+n,2)$ and for any $n\in\N$
one has
\begin{gather}\label{eq:epsilon:0:2corr}
\epsilon\big(0^{n-1},d_1,d_2\big)\ge -\frac{2}{6g-1} .
\end{gather}
\end{Corollary}
\begin{proof}If one of $d_1$, $d_2$
is equal to zero, the statement for arbitrary $n$
follows from~\eqref{eq:epsilon:0:1corr}, so from now on
we assume
that both $d_1$, $d_2$ are strictly positive.
For $n=1$ the statement follows directly from~\eqref{eq:main:bounds}.
This serves us as a base of induction
in $n$. Suppose that for all $n=1,\dots,k$ the statement
is true. Let us prove it for $n=k+1$:
\begin{gather*}
\big\langle \tau_0^{k+1}\tau_{d_1}\tau_{d_2}\big\rangle_{g,k+3}
=\big\langle \tau_0^k\tau_{d_1-1}\tau_{d_2}\big\rangle_{g,k+2}
+\big\langle \tau_0^k\tau_{d_1}\tau_{d_2-1}\big\rangle_{g,k+2}\\
\hphantom{\big\langle \tau_0^{k+1}\tau_{d_1}\tau_{d_2}\big\rangle_{g,k+3}}{} \ge
\left(1-\frac{2}{6g-1}\right)\cdot
\big(\big\lfloor \tau_0^k\tau_{d_1-1}\tau_{d_2}\big\rfloor_{g,k+2}
+\big\lfloor \tau_0^k\tau_{d_1}\tau_{d_2-1}\big\rfloor_{g,k+2}
\big)\\
\hphantom{\big\langle \tau_0^{k+1}\tau_{d_1}\tau_{d_2}\big\rangle_{g,k+3}}{}
 =
\left(1-\frac{2}{6g-1}\right)
\cdot\big(1+\delta_{\rm string}\big(0^{k+1},d_1,d_2\big)\big)\cdot
\big\lfloor \tau_0^{k+1}\tau_{d_1}\tau_{d_2}\big\rfloor_{g,k+3}\\
\hphantom{\big\langle \tau_0^{k+1}\tau_{d_1}\tau_{d_2}\big\rangle_{g,k+3}}{} \ge
\left(1-\frac{2}{6g-1}\right)\cdot
\big\lfloor \tau_0^{k+1}\tau_{d_1}\tau_{d_2}\big\rfloor_{g,k+3} ,
\end{gather*}
where the first equality is the string equation;
inequality between the first and the second lines is
the assumption of the induction; the equality
between the second and the third line is
equation~\eqref{eq:string:floor}; the inequality
between the third and the forth line is
an implication of~\eqref{eq:delta:string:is:positive}
and of the fact that all the factors in both lines are
positive.
\end{proof}

\begin{Corollary}\label{cor:L:equals:0}
For any $g,n\in\N$ and for any
$\boldsymbol{d}\in\Pi_0(3g-3+n,n)$ one has
\begin{gather}\label{eq:epsilon:is:smaller:than:lambda:0}
1+\epsilon(\boldsymbol{d})\ge\lambda(g,0) .
\end{gather}
\end{Corollary}

\begin{proof}Recalling convention~\eqref{eq:lambda:0} for $\lambda(g,0)$ we conclude that
for $n=1$ inequality~\eqref{eq:epsilon:is:smaller:than:lambda:0}
follows from~\eqref{eq:epsilon:for:1:correlators};
for $n=2$ inequality~\eqref{eq:epsilon:is:smaller:than:lambda:0}
follows from~\eqref{eq:main:bounds};
for $n\ge 3$ inequality~\eqref{eq:epsilon:is:smaller:than:lambda:0}
corresponds to~\eqref{eq:epsilon:0:2corr}.
\end{proof}

\begin{Lemma}\label{lm:dilaton}
Let $\boldsymbol{d}\in\Pi(3g-3+n,n)$.
Define $\delta_{\rm dilaton}(1,\boldsymbol{d})$
by equation
\begin{gather}\label{eq:dilaton:floor}
\lfloor \tau_1\tau_{d_1} \cdots \tau_{d_n}\rfloor_{g,n+1}
\cdot\big(1+\delta_{\rm dilaton}(1,\boldsymbol{d})\big)
=(2g-2+n)
\lfloor \tau_{d_1} \cdots \tau_{d_n}\rfloor_{g,n} .
\end{gather}
Then
\[
\delta_{\rm dilaton}(1,\boldsymbol{d})
=\frac{n-3}{6g-3+2n} .
\]
In particular,
\begin{gather}\label{eq:delta:dilaton}
\delta_{\rm dilaton}(1,\boldsymbol{d})
\begin{cases}
\ge 0&\text{when }n\ge 3 ,
\\
= -\dfrac{1}{6g+1}
&\text{when }n=2 ,
\vspace{1mm}\\
= -\dfrac{2}{6g-1}
&\text{when }n=1 .
\end{cases}
\end{gather}
\end{Lemma}

\begin{proof}Dividing both sides of equation~\eqref{eq:dilaton:floor} by
$\lfloor \tau_1\tau_{d_1} \cdots \tau_{d_n}\rfloor_{g,n+1}$,
applying definition~\eqref{eq:floor:brackets}
and canceling common factors in the numerator and in the denominator
of the resulting expression we get
\[
1+\delta_{\rm dilaton}(1,\boldsymbol{d})
=\frac{3(2g-2+n)}{6g-3+2n}
=\frac{6g-6+3n}{6g-3+2n} .\tag*{\qed}
\]\renewcommand{\qed}{}
\end{proof}

\begin{Corollary}\label{cor:epsilon:L:1}
For any partition
$\boldsymbol{d}\in\Pi_1(3g-3+n,n)$ one has
\begin{gather}\label{eq:epsilon:L:1}
1+\epsilon(\boldsymbol{d})\ge
\left(1-\frac{1}{6g+1}\right)
\cdot\left(1-\frac{2}{6g-1}\right) .
\end{gather}
\end{Corollary}

\begin{proof}We have seen in~\eqref{eq:epsilon:for:1:correlators}
that for all $1$-correlators we have
$\epsilon(3g-2)=0$,
so for $n=1$ the statement is true.
For $n=2$ the statement is a direct implication
of equation~\eqref{eq:main:bounds}.
Suppose that $n\ge 3$.

By Remark~\ref{rm:permutation:of:d}, the quantity
$\epsilon(\boldsymbol{d})$ does not change under any permutation of the entries of $\boldsymbol{d}$. Thus, we can permute the first $n-2$ elements of the partition without affecting the value of $\epsilon(\boldsymbol{d})$, in particular, we can place them in the growing order. Since the sum of the first $n-2$
elements is less than or equal to $1$ either they are all
equal to $0$ or they form the sequence $(0,\dots,0,1)$
after such reordering. If they all are equal to $0$, the
statement follows from equation~\eqref{eq:epsilon:0:2corr}
from Corollary~\eqref{cor:epsilon:0:2corr}.

It remains to consider the case when $n\ge 3$ and when the first $n-2$ elements form a sequence $(0,\dots,0,1)$. We prove first the desired inequality for partitions of the form $(1,d_1,d_2)$:
\begin{gather*}
\langle \tau_1\tau_{d_1}\tau_{d_2}\rangle_{g,3}
=(2g)\cdot\langle \tau_{d_1} \tau_{d_2}\rangle_{g,2}
\ge\left(1-\frac{2}{6g-1}\right)\cdot
(2g)\lfloor \tau_{d_1}\tau_{d_2}\rfloor_{g,2}
\\
\hphantom{\langle \tau_1\tau_{d_1}\tau_{d_2}\rangle_{g,3}}{}
=\left(1-\frac{2}{6g-1}\right)\cdot
\lfloor \tau_1\tau_{d_1}\tau_{d_2}\rfloor_{g,3}
\cdot\big(1+\delta_{\rm dilaton}(1,d_1,d_2)\big)
\\
\hphantom{\langle \tau_1\tau_{d_1}\tau_{d_2}\rangle_{g,3}}{}
\ge\left(1-\frac{2}{6g-1}\right)\left(1-\frac{1}{6g+1}\right)
\cdot\lfloor \tau_1\tau_{d_1}\tau_{d_2}\rfloor_{g,3} .
\end{gather*}
Here the first equality is the dilaton equation;
the inequality which follows is
equation~\eqref{eq:main:bounds}; the equality
between the first two lines is equation~\eqref{eq:dilaton:floor}
and the inequality between the second and the third line
is based on equation~\eqref{eq:delta:dilaton}.

To complete the proof of Corollary~\ref{cor:epsilon:L:1} we prove it for partitions of the form
$\big(0^{k+1},1,d_1,d_2\big)$ by induction in $k\ge 0$. The proof follows line-by-line the proof of Corollary~\ref{cor:epsilon:0:2corr}.
\end{proof}

\begin{Lemma}\label{lm:Virasoro}
Let $\boldsymbol{d}\in\Pi(3g-3+n-k,n)$, where $k,n\in\N$.
Define $\delta_{\rm Virasoro}(k+1,\boldsymbol{d})$
by equation
\begin{gather}
\lfloor\tau_{k+1}\tau_{d_1}\cdots\tau_{d_n}\rfloor_g
\big(1+\delta_{\rm Virasoro}(k+1,\boldsymbol{d})\big)
\nonumber\\
\qquad{} =\frac{1}{(2k+3)!!}
\Bigg(\sum_{j=1}^n
\frac{(2k+2d_j+1)!!}{(2d_j-1)!!}\lfloor\tau_{d_1}\cdots
\tau_{d_{j}+k}\cdots\tau_{d_n}\rfloor_g\nonumber\\
\qquad\quad{} +\frac{1}{2}\sum_{\substack{r+s=k-1\\r,s\ge 0}}
(2r+1)!!(2s+1)!!\lfloor\tau_r\tau_s\tau_{d_1}\cdots\tau_{d_n}\rfloor_{g-1}\Bigg) .\label{eq:Virasoro:floor}
\end{gather}
Then
\begin{gather}\label{eq:delta:Virasoro}
\delta_{\rm Virasoro}(k+1,\boldsymbol{d})
=\frac{n-3}{6g-3+2n}-\frac{2k(2n-5)}{(6g-3+2n)(6g-5+2n)} .
\end{gather}
\end{Lemma}

\begin{proof}Dividing both sides of equation~\eqref{eq:Virasoro:floor} by
$\lfloor \tau_{k+1}\tau_{d_1} \cdots \tau_{d_n}\rfloor_{g,n+1}$,
applying definition~\eqref{eq:floor:brackets}
and canceling common factors in the numerator and in the denominator
of the resulting expression we get
\begin{gather*}
1+\delta_{\rm Virasoro}(k+1,\boldsymbol{d})\\
\qquad{} =
\frac{1}{6g-3+2n}\cdot\left(
\big((2d_1+1)+\dots+(2d_n+1)\big)
+\frac{1}{2}\cdot k\cdot \frac{24g}{6g-5+2n}\right) \\
 \qquad{} =\frac{1}{6g-3+2n}\cdot\left(
\big(6g-6+3n-2k\big)+2k-2k\cdot\frac{2n-5}{6g-5+2n}\right)\\
\qquad{} = 1+\frac{n-3}{6g-3+2n}-\frac{2k(2n-5)}{(6g-3+2n)(6g-5+2n)} .\tag*{\qed}
\end{gather*}\renewcommand{\qed}{}
\end{proof}

\begin{Remark}\label{rm:lower:bound:Virasoro}
In expression~\eqref{eq:Virasoro:floor} we ignored the third term in the Virasoro constraints. Since this third term is, clearly, positive,
this is suitable for getting a lower bound
instead of exact asymptotics.
It is widely believed that the third term
of Virasoro constraints becomes negligible in large genera. We expect that technique from~\cite{Aggarwal:Siegel:Veech} might be useful for replacing
the lower bound in~\eqref{eq:main:lower:bound} by the exact asymptotics
under strengthening restrictions on $\alpha$.
\end{Remark}

We shall need the following technical corollary of Lemma~\ref{lm:Virasoro}.

\begin{Corollary}\label{cor:delat:Virasoro:bound}
Let $k,n$ be integers satisfying $k\ge 0$, $n\ge 2$.
Let $\boldsymbol{d}$ be a partition $\boldsymbol{d}\in\Pi(3g-3+n-k,n)$, such that
$k+1\le d_j$ for $j=1,\dots,n-2$,
and $k+d_1+\dots+d_{n-2}\le\frac{3}{2}g$.
Then
\begin{gather}\label{eq:delat:Virasoro:bound}
\delta_{\rm Virasoro}(k+1,\boldsymbol{d})\ge -\frac{1}{6g+1} .
\end{gather}
\end{Corollary}

\begin{proof}Use expression~\eqref{eq:delta:Virasoro} for
$\delta_{\rm Virasoro}(k+1,\boldsymbol{d})$.
When $n=2$ we get
\[
\delta_{\rm Virasoro}(k+1,\boldsymbol{d})=-\frac{1}{6g+1}+\frac{2k}{(6g+1)(6g-1)}\ge-\frac{1}{6g+1} .
\]

Let $n\ge 3$. By assumption, $(k+1)\le d_j$ for $j=1,\dots,d_{n-2}$, so $(k+1)(n-2)\le \frac{3}{2}g$, and hence $2k(2n-5)< 6g$, which implies that
\begin{gather*}
\delta_{\rm Virasoro}((k+1,\boldsymbol{d}))=\frac{n-3}{6g-3+2n}-\frac{2k(2n-5)}{(6g-3+2n)(6g-5+2n)}\\
\hphantom{\delta_{\rm Virasoro}((k+1,\boldsymbol{d}))}{} \ge -\frac{2k(2n-5)}{(6g-3+2n)(6g-5+2n)}
\ge -\frac{6g}{(6g+3)(6g+1)}>-\frac{1}{6g+1} .\!\!\!\tag*{\qed}
\end{gather*}\renewcommand{\qed}{}
\end{proof}

Before passing to the step of induction, we recapitulate the properties of the function $\lambda(g,L)$ defined in~\eqref{eq:lambda}.
Recall that the arguments $g$, $L$ of $\lambda(g,L)$ are nonnegative integers
satisfying $g>L$. Inequalities~\eqref{eq:lambda:inequalities}--\eqref{eq:lambda:inequality:ters}
below follow from the definition of $\lambda(g,L)$. Each inequality is
applied to those ordered pairs $g$, $L$ for which the argument of $\lambda$ on
both sides of the inequality belongs to the domain of definition of~$\lambda$. We have
\begin{gather}\label{eq:lambda:inequalities}
1>\lambda(g+1,L)>\lambda(g,L)>\lambda(g,L+1)>0 ,
\end{gather}
and
\begin{gather}\label{eq:lambda:inequality:bis}
\left(1-\frac{1}{6g+1}\right)
\cdot\lambda(g-1,L-1)=\lambda(g,L) .
\end{gather}
Combining the latter two relations we get
\begin{gather}
\left(1-\frac{1}{6g+1}\right)\cdot\lambda(g,L)>\left(1-\frac{1}{6(g-1)+1}\right)\cdot\lambda(g,L)\nonumber\\
\hphantom{\left(1-\frac{1}{6g+1}\right)\cdot\lambda(g,L)}{} >\left(1-\frac{1}{6(g-1)+1}\right)
\cdot\lambda(g-1,L) =\lambda(g,L+1) .\label{eq:lambda:inequality:ters}
\end{gather}

\begin{Proposition}[step of induction] \label{eq:prop:step:of:induction}
Suppose that for some nonnegative integer $L_0$
the following uniform bound is valid:
for all integers $g$, $L$ satisfying $g> L_0$, $0\le L\le L_0$,
for all partitions $\boldsymbol{d}\in\Pi_L(3g-2+n,n+1)$,
where $n\ge 0$, one has
\[
1+\epsilon(\boldsymbol{d})\ge \lambda(g,L) .
\]
Then for all integers $g$, $L$
satisfying $g> L_0+1$,
$0\le L\le L_0+1$,
for all partitions $\boldsymbol{d}\in\Pi_L(3g-2+n,n+1)$,
where $n\ge 0$, one also has
\[
1+\epsilon(\boldsymbol{d})\ge \lambda(g,L) .
\]
\end{Proposition}

\begin{proof}
We warn the reader that the total number of elements of the
partition is denoted in
Proposition~\ref{eq:prop:step:of:induction} by $n+1$ and
not by $n$ as in Theorem~\ref{th:Main:Theorem}. This allows
us to use formulae from the key Lemmas~\ref{lm:dilaton}
and~\ref{lm:Virasoro} without adjustments.

By convention $\Pi_L(3g-2,1)=\Pi(3g-2,1)$ and
$\Pi_L(3g-1,2)=\Pi(3g-1,2)$ for any $L\in\Z_{\ge 0}$.
We have seen in~\eqref{eq:epsilon:for:1:correlators}
that for all $1$-correlators we have
$\epsilon(3g-2)=0$,
so for $n=0$ the statement is trivially true.
For $n=1$ the statement is a direct implication
of inequality~\eqref{eq:main:bounds}:
\[
1+\epsilon(\boldsymbol{d})
\ge \left(1-\frac{2}{6g-1}\right)
=\lambda(g,0)
\ge\lambda(g,L) .
\]
Thus, from now on we can assume that $n\ge 2$. Let $\boldsymbol{d}\in\Pi_L(3g-2+n,n+1)$. If $L\le L_0$, then the statement makes part of the induction assumption. Hence, from now on we can assume that
\[
\boldsymbol{d}\in\Pi_{L_0+1}(3g-2+n,n+1)\setminus\Pi_{L_0}(3g-2+n,n+1) ,
\]
where $n\ge 2$ and $g> L_0+1$. This implies that
$d_1+\dots+d_{n-1}=L_0+1$ and, hence,
\begin{gather*}%\label{eq:sum:of:last:two}
d_{n}+d_{n+1}=3g-2+n-(L_0+1)> 2L_0+2 .
\end{gather*}

By Remark~\ref{rm:permutation:of:d}, the quantity $\epsilon(\boldsymbol{d})$ does not change under any permutation of the entries of~$\boldsymbol{d}$. Place to the leftmost position the smallest strictly positive element among the first $n-1$ elements. This operation does not change the last two elements of the partition and does not change the sum of its first $n-1$ elements. Denote the resulting partition by $(k+1,d_1,\dots,d_n)$. To prove the proposition we have to prove the inequality
\begin{gather}\label{eq:desired:inequality}
1+\epsilon(k+1,d_1,\dots,d_n)\ge \lambda(g,L_0+1) ,
\end{gather}
where
\begin{gather}
k+d_1+d_2+\dots+d_{n-2}=L_0 ,\qquad n\ge 2,\qquad k\ge 0 ,\nonumber\\
 k+1\le\min_{\substack{1\le i\le n-2\\d_i>0}} d_i ,\qquad g>L_0+1\ge 1 .\label{eq:restrictions}
\end{gather}

We consider the special case $k=0$ separately. In this special case,
when $n=2$ the
partition $(k+1,d_1,\dots,d_n)$ becomes $(1,d_1,d_2)$ and the
desired inequality is proved in~\eqref{eq:epsilon:L:1}. Assume that
$k=0$ and $n\ge 3$. By~\eqref{eq:delta:dilaton} we
have
\begin{gather}\label{eq:delta:dilaton:positive}
\delta_{\rm dilaton}(1,d_1,\dots,d_n)\ge 0 ,\qquad\text{for }n\ge 3 .
\end{gather}
Thus, for any $g> L_0+1$ we have
\begin{gather*}
\big(1+\epsilon(1,d_1,\dots,d_n)\big)\cdot \lfloor \tau_1\tau_{d_1} \cdots \tau_{d_n}\rfloor_{g,n+1}=
\langle \tau_1\tau_{d_1} \cdots \tau_{d_n}\rangle_{g,n+1}\\
\qquad{} =(2g-2+n)\langle \tau_{d_1} \cdots \tau_{d_n}\rangle_{g,n}
\ge (2g-2+n)\cdot\big(\lambda(g,L_0)\cdot
\lfloor \tau_{d_1} \cdots \tau_{d_n}\rfloor_{g,n}\big)
\\
\qquad{} =\lambda(g,L_0)\cdot
\lfloor \tau_1\tau_{d_1} \cdots \tau_{d_n}\rfloor_{g,n+1}
\cdot\big(1+\delta_{\rm dilaton}(1,\boldsymbol{d})\big)
\\
\qquad{} \ge \lambda(g,L_0)\cdot
\lfloor \tau_1\tau_{d_1} \cdots \tau_{d_n}\rfloor_{g,n+1}
> \lambda(g,L_0+1)
\cdot
\lfloor \tau_1\tau_{d_1} \cdots \tau_{d_n}\rfloor_{g,n+1} .
\end{gather*}
Here the first equality is
definition~\eqref{eq:def:of:epsilon}
of $\epsilon(1,d_1,\dots,d_n)$; the second equality
is the string equation~\eqref{eq:string}; the
inequality in the middle of the second line is the
induction assumption; the equality in the beginning of the
third line is definition~\eqref{eq:dilaton:floor}
of $(1+\delta_{\rm dilaton}(1,\boldsymbol{d}))$;
the inequality in the beginning of the last line
is a direct implication of~\eqref{eq:delta:dilaton:positive};
the last inequality is an implication of~\eqref{eq:lambda:inequalities}.

Suppose now that $k\ge 1$. We first prove the desired inequality~\eqref{eq:desired:inequality} in the special case when
\begin{gather}\label{eq:all:dj:strictly:positive}
d_j\ge 1 ,\qquad\text{for }j=1,\dots,n-2
\end{gather}
and then prove it in the most general situation when some of $d_j$ (possibly all of them) are equal to $0$.

Note that by assumption, $k+1$ is less than or equal
to any strictly positive element among $d_1,\dots,d_{n-2}$,
so inequalities~\eqref{eq:all:dj:strictly:positive},
actually, imply that
\begin{gather*}%\label{eq:all:dj:strictly:positive:bis}
d_j\ge k+1 ,\qquad\text{for }j=1,\dots,n-2.
\end{gather*}
Also, from~\eqref{eq:restrictions} we get
\[
k+d_1+d_2+\dots+d_{n-2}=L_0<\frac{3g}{2} .
\]
Thus, the partition $(k+1,d_1,\dots,d_n)$
satisfies assumptions of Corollary~\ref{cor:delat:Virasoro:bound}.

Let $g> L_0+1$. Under the above assumptions we have
\begin{gather*}
\big(1+\epsilon(k+1,d_1,\dots,d_n)\big)\cdot\lfloor\tau_{k+1}\tau_{d_1}\cdots\tau_{d_n}\rfloor_g
=\langle\tau_{k+1}\tau_{d_1}\cdots\tau_{d_n}\rangle_g
 \\
 \qquad{} \ge
\frac{1}{(2k+3)!!}
\Bigg(\sum_{j=1}^n
\frac{(2k+2d_j+1)!!}{(2d_j-1)!!}\langle\tau_{d_1}\cdots
\tau_{d_{j}+k}\cdots\tau_{d_n}\rangle_g\\
\qquad\quad{} +\frac{1}{2}\sum_{\substack{r+s=k-1\\r,s\ge 0}}
(2r+1)!!(2s+1)!!\langle\tau_r\tau_s\tau_{d_1}\cdots\tau_{d_n}\rangle_{g-1}
\Bigg)
\\ \qquad{} \ge
\frac{1}{(2k+3)!!}
\Bigg(\sum_{j=1}^{n-2}
\frac{(2k+2d_j+1)!!}{(2d_j-1)!!}
\cdot\lambda(g,L_0)\cdot
\lfloor\tau_{d_1}\cdots
\tau_{d_{j}+k}\cdots\tau_{d_n}\rfloor_g
\\
\qquad\quad {}+\frac{(2k+2d_{n-1}+1)!!}{(2d_{n-1}-1)!!}
\cdot\lambda(g,L_0-k)\cdot
\lfloor\tau_{d_1}\cdots\tau_{d_{n-2}}
\tau_{d_{n-1}+k}\tau_{d_n}\rfloor_g
\\
\qquad\quad {}+\frac{(2k+2d_{n}+1)!!}{(2d_{n}-1)!!}
\cdot\lambda(g,L_0-k)\cdot
\lfloor\tau_{d_1}\cdots\tau_{d_{n-1}}
\tau_{d_{n}+k}\rfloor_g
\\
\qquad\quad {}+\frac{1}{2}\sum_{\substack{r+s=k-1\\r,s\ge 0}}
(2r+1)!!(2s+1)!!
\cdot\lambda(g-1,L_0-1)\cdot
\lfloor\tau_r\tau_s\tau_{d_1}\cdots\tau_{d_n}\rfloor_{g-1}
\Bigg)
\\
\qquad{} \ge\lambda(g,L_0)\cdot\lfloor\tau_{k+1}\tau_{d_1}\cdots\tau_{d_n}\rfloor_g
\big(1+\delta_{\rm Virasoro}(k+1,\boldsymbol{d})\big)
\\
\qquad{} \ge\left(1-\frac{1}{6g+1}\right)\cdot\lambda(g,L_0)\cdot\lfloor\tau_{k+1}\tau_{d_1}\cdots\tau_{d_n}\rfloor_g
\\
\qquad{} \ge\lambda(g,L_0+1)\cdot\lfloor\tau_{k+1}\tau_{d_1}\cdots\tau_{d_n}\rfloor_g .
\end{gather*}
Here the first equality is
the definition~\eqref{eq:def:of:epsilon}
of $\epsilon(1,d_1,\dots,d_n)$;
the first inequality
is an instant corollary of the
Virasoro constraints in which we omitted
the terms in the third line of~\eqref{eq:virasoro}.
The second inequality is the induction assumption.
The third inequality combines
the inequality $\lambda(g,L_0-k)>\lambda(g,L_0)$
which follows from~\eqref{eq:lambda:inequalities},
the inequality $\lambda(g-1,L_0-1)>\lambda(g,L_0)$
which follows from~\eqref{eq:lambda:inequality:bis},
and the definition~\eqref{eq:Virasoro:floor}
of $\delta_{\rm Virasoro}(k+1,\boldsymbol{d})$.
The inequality
$\big(1+\delta_{\rm Virasoro}(k+1,\boldsymbol{d})\big)\ge
\big(1-\frac{1}{6g+1}\big)$ is justified by~\eqref{eq:delat:Virasoro:bound}.
The last inequality is justified in~\eqref{eq:lambda:inequality:ters}.

It remains to prove inequality~\eqref{eq:desired:inequality},
without extra assumptions~\eqref{eq:all:dj:strictly:positive}.
In other words, we have to prove the inequality
\[
1+\epsilon\big(0^s,k+1,d_1,\dots,d_{n-s}\big)\ge \lambda(g,L_0+1) .
\]

The case $n-s=0$ follows from~\eqref{eq:epsilon:0:1corr}. For $n-s=1$ inequality~\eqref{eq:epsilon:0:2corr} implies
\[
1+\epsilon\big(0^s,k+1,d_1\big)\ge 1-\frac{2}{6g-1} =\lambda(g,0) \ge \lambda(g,L_0+1) .
\]
Thus, we may assume that $n-s\ge 2$ and that the following inequalities are valid:
\[
\begin{cases}
k+d_1+d_2+\dots+d_{n-s-2}=L_0 ,
\\ n\ge s\ge 0 ,
\\ k\ge 1 ,
\\ d_j> k\quad\text{for }j=1,\dots,n-s-2 .
\end{cases}
\]

We proceed by induction in $s$. For $s=0$, which serves us as a base of induction, the statement is already proved. We
perform a step of induction as follows
\begin{gather*}
\big(1+\epsilon\big(0^{s+1},k+1,d_1,\dots,d_{n-s}\big)\big)
\cdot\big\lfloor\tau_0^{s+1}\tau_{k+1}\tau_{d_1}\cdots\tau_{d_{n-s}}\big\rfloor_g
 =\big\langle\tau_0^{s+1}\tau_{k+1}\tau_{d_1}\cdots\tau_{d_{n-s}}\big\rangle_g
\\
\qquad{} =\big\langle\tau_0^s\tau_k\tau_{d_1}\cdots\tau_{d_{n-s}}\big\rangle_g
+\big\langle\tau_0^s\tau_{k+1}\tau_{d_1-1}\cdots\tau_{d_{n-s}}\big\rangle_g
+\dots
+\big\langle\tau_0^s\tau_{k+1}\tau_{d_1}\cdots\tau_{d_{n-s}-1}\big\rangle_g
\\
\qquad{} =\big(1+\epsilon\big(0^s,k,d_1,\dots,d_{n-s}\big)\big)
\cdot\big\lfloor\tau_0^s\tau_k\tau_{d_1}\cdots\tau_{d_{n-s}}\big\rfloor_g
\\
\qquad\quad{}+\big(1+\epsilon\big(0^s,k+1,d_1-1,\dots,d_{n-s}\big)\big)
\cdot\big\lfloor\tau_0^s\tau_{k+1}\tau_{d_1-1}\cdots\tau_{d_{n-s}}\big\rfloor_g
+\cdots
\\
\qquad\quad{} +\big(1+\epsilon\big(0^s,k+1,d_1,\dots,d_{n-s}-1\big)\big)
\cdot\big\lfloor\tau_0^s\tau_{k+1}\tau_{d_1}\cdots\tau_{d_{n-s}-1}\big\rfloor_g
\\
\qquad{} \ge\lambda(g,L_0)
\cdot\big\lfloor\tau_0^s\tau_k\tau_{d_1}\cdots\tau_{d_{n-s}}\big\rfloor_g
+\lambda(g,L_0)
\cdot\big\lfloor\tau_0^s\tau_{k+1}\tau_{d_1-1}\cdots\tau_{d_{n-s}}\big\rfloor_g+\cdots
\\
\qquad\quad{} +\lambda(g,L_0)
\cdot\big\lfloor\tau_0^s\tau_{k+1}\tau_{d_1}\cdots\tau_{d_{n-s-2}-1}\tau_{d_{n-s-1}}\tau_{d_{n-s}}\big\rfloor_g\\
\qquad\quad{} +\lambda(g,L_0+1)
\cdot\big\lfloor\tau_0^s\tau_{k+1}\tau_{d_1}\cdots\tau_{d_{n-s-2}}\tau_{d_{n-s-1}-1}\tau_{d_{n-s}}\big\rfloor_g
\\
\qquad\quad{}
+\lambda(g,L_0+1)
\cdot\big\lfloor\tau_0^s\tau_{k+1}\tau_{d_1}\cdots\tau_{d_{n-s-2}}\tau_{d_{n-s-1}}\tau_{d_{n-s}-1}\big\rfloor_g
\\
\qquad{} \ge\lambda(g,L_0+1)\big(
\big\lfloor\tau_0^s\tau_k\tau_{d_1}\cdots\tau_{d_{n-s}}\big\rfloor_g
+\big\lfloor\tau_0^s\tau_{k+1}\tau_{d_1-1}\cdots\tau_{d_{n-s}}\big\rfloor_g
+\cdots\\
\qquad\quad{}
+\big\lfloor\tau_0^s\tau_{k+1}\tau_{d_1}\cdots\tau_{d_{n-s}-1}\big\rfloor_g
\big)\\
\qquad{} =
\lambda(g,L_0+1)
\cdot\big\lfloor\tau_0^{s+1}\tau_{k+1}\tau_{d_1}\cdots\tau_{d_{n-s}}\big\rfloor_g
\cdot\big(1+\delta_{\rm string}\big(0^{s+1},k+1,d_1,\dots,d_{n-s}\big)\big)\\
\qquad{} \ge \lambda(g,L_0+1)
\cdot\big\lfloor\tau_0^{s+1}\tau_{k+1}\tau_{d_1}\cdots\tau_{d_{n-s}}\big\rfloor_g.
\end{gather*}
Here the first equality is the definition~\eqref{eq:def:of:epsilon} of $\epsilon\big(0^{s+1},k+1,d_1,\dots,d_{n-s}\big)$. The second equality is the string equation~\eqref{eq:string}.
(Recall that by convention, if one of $d_{n-s-1}$ or $d_{n-s}$ is equal to zero, the term, containing the negative index $d_{n-s-1}-1$ or $d_{n-s}-1$ respectively, is missing in the string equation and below.)
The equality which follows, is equation~\eqref{eq:def:of:epsilon} applied to every term of the resulting expression. The inequality, where $\lambda$ appears on the left-hand side for the first time, is the induction assumption applied to each term. The next inequality follows from the inequality $\lambda(g,L_0)>\lambda(g,L_0+1)$,
see~\eqref{eq:lambda:inequalities}. The equality which follows is the definition~\eqref{eq:string:floor} of $\delta_{\rm string}\big(0^{s+1},k+1,d_1,\dots,d_{n-s}\big)$. The last inequality is justified by~\eqref{eq:delta:string:is:positive}.
\end{proof}

\begin{proof}[Proof of Theorem~\ref{th:Main:Theorem}]
For $L=0$ Theorem~\ref{th:Main:Theorem} follows from
Corollary~\ref{cor:L:equals:0}. For $L=1$ it follows from~\eqref{eq:epsilon:L:1}
combined with the fact that $\big(1-\frac{1}{6g+1}\big)
\cdot\big(1-\frac{2}{6g-1}\big)>\lambda(g,1)$.
For $L>1$ we apply recursively Proposition~\ref{eq:prop:step:of:induction}.
\end{proof}
